\section{Introduction}
\label{sec:intro}

An argumentation task asks for a position and a case for it. A debater, a policy analyst, or an advocate must take a stance, make claims that support it, back each claim with evidence, and answer the other side. Deep research pipelines such as OpenAI Deep Research \citep{openai2025deepresearch} and STORM \citep{shao2024storm} search the web, read what they find, and write a cited report, but their objective is a summary. Asked whether the Federal Reserve should keep interest rates high, they describe both views and stop. For a survey that is the right behavior. For a debate it is a non-answer.

The two kinds of task differ in how research should be organized. A summary can grow bottom-up: search, read, search again where the text is thin, and write up what was found. An argument has a shape before the research starts. The root claim rests on a few supporting claims, each of which rests on evidence, and the research has to fill that structure. This is decomposition, the central move of computational thinking, applied to the logic of an argument instead of the structure of a program: a large question is split into parts that can be answered on their own and then recombined.

\sys{} is a deep research pipeline built on that idea. It first plans: a language model, given the motion and a short literature search, generates a directed acyclic graph (DAG) of sub-questions, in which every non-leaf node also carries instructions for composing its answer from its children. It then researches: leaves are answered by web search, parents are composed from their children in topological order, and a parent whose composed answer is incomplete spawns a sub-tree to fill it. It finally writes: the answered DAG becomes an outline and then a report that argues the root answer. \Cref{fig:pipeline} shows the three phases.

The pipeline was compared with OpenAI Deep Research and STORM in two original evaluations, six automated debates and ten reports scored blind by one rater, in which it scored highest on argument quality and was far more decisive. A second look at the retained reports, which needs no rater, refines both results. Paired lexical measurements confirm that \sys{} commits where STORM hedges, but show that its conclusions are often softer than its titles and that it cites densely from few sources. A citation attribution test then finds a problem the original rubric could not see: the citations in the evaluated reports do not match their claims better than chance. The cause is a provenance failure in recursive composition, in which citation numbers lose their meaning as answers move up the graph, and it can be traced to two lines of the pipeline's design.

\Cref{sec:problem} defines the task and the property of provenance that the pipeline must preserve, \Cref{sec:method} describes the pipeline, \Cref{sec:experiments,sec:results} give the original evaluation with its limits, \Cref{sec:analysis} analyzes what the reports contain, \Cref{sec:citations} presents the attribution test, the diagnosis, and the fix, and \Cref{sec:discussion} states what the results do and do not support. The diagnosis applies to any system that composes cited answers recursively.
